% TOP_PROBLEM: {"id": 1485, "title": "Volume Conjecture", "category_id": 7, "category": {"id": 7, "name": "topology", "display_name": "Topology", "description": "Properties preserved under continuous deformations.", "slug": "topology", "order_index": 7, "created_at": "2026-07-31T15:26:25.671Z"}, "status": "open"}
% FIRST_POSTED: null
% ENTRY_KIND: statement_only
% Imported from: batch08.tex; UnsolvedMath ID 1485
\documentclass[11pt]{article}
\usepackage[margin=25mm]{geometry}
\usepackage{fontspec}
\setmainfont{Latin Modern Roman}
\usepackage{amsmath,amssymb,mathtools}
\usepackage{unicode-math}
\setmathfont{Latin Modern Math}
\usepackage{xurl,hyperref}
\hypersetup{colorlinks=true,urlcolor=blue}
\setcounter{secnumdepth}{0}
\setlength{\parindent}{0pt}
\setlength{\parskip}{5pt}
\setlength{\emergencystretch}{3em}
\newcommand{\sref}[2]{\hyperlink{source:#1:#2}{[S#2]}}
\newcommand{\cataloguescope}[1]{\paragraph{Recorded target.}#1}
\begin{document}
% UnsolvedMath ID 1485; source catalogue code TOP-007
\setcounter{section}{1484}
\section{Volume Conjecture}\label{1485}
{\small\sffamily UnsolvedMath ID 1485 · Topology}
\subsection{Definitions and mathematical statement}

\paragraph{Shared notation.}$\mathbb N=\{1,2,\ldots\}$, and $\mathbb Z,\mathbb Q,\mathbb R,\mathbb C$ denote the integers, rationals, reals and complex numbers. Any other conventions are specified locally.


Let \(K\subset S^3\) be a tame knot, and let \(J_N(K;q)\)
be its \(N\)-colored Jones polynomial, defined using the
\(N\)-dimensional irreducible \(\mathfrak{sl}_2\) representation,
with zero framing and normalization \(J_N(\text{unknot};q)=1\).
Use the convention \(q=e^{2\pi i/N}\); in a convention using
\(s=q^{1/2}\), this is \(s=e^{\pi i/N}\).
For a hyperbolic knot, \(\operatorname{Vol}(S^3\setminus K)\) is
the volume of the complete finite-volume hyperbolic metric on its
complement. More generally, cut the compact knot exterior along its
canonical essential tori, and let \(V(K)\) be the sum of the
finite hyperbolic volumes of the hyperbolic pieces (zero if there are none).
Equivalently \(V(K)=v_3\|K\|\), where \(v_3\) is the volume of a
regular ideal hyperbolic tetrahedron and \(\|K\|\) is the simplicial
volume of the exterior. In particular \(V(K)=\operatorname{Vol}(S^3\setminus K)\)
for hyperbolic knots.
\paragraph{Statement recorded in the supplied source.}
For every hyperbolic knot \(K\), is
\[
 \lim_{N\to\infty}\frac{2\pi}{N}
      \log\bigl|J_N(K;e^{2\pi i/N})\bigr|
       =\operatorname{Vol}(S^3\setminus K)?
\]
The standard all-knot extension proposed in the cited source asks,
for every knot \(K\), whether the same limit exists and equals \(V(K)\).
When a Jones value is zero, its logarithm is interpreted as \(-\infty\). \sref{1485}{2}
\subsection{Short English statement}
Does the exponential growth of normalized colored Jones values at successive roots of unity recover hyperbolic knot-complement volume? The all-knot extension asks for the sum of the hyperbolic volumes of the exterior's pieces.
\subsection{Sources}
\begin{description}
\item[\hypertarget{source:1485:1}{S1}] ULAM.AI and the UnsolvedMath Contributors, \emph{UnsolvedMath: A Curated Collection of Open Mathematics Problems}, record 1485 (TOP-007). CC BY 4.0. \url{https://huggingface.co/datasets/ulamai/UnsolvedMath}.
\item[\hypertarget{source:1485:2}{S2}] Hitoshi Murakami and Jun Murakami, \emph{The Colored Jones Polynomials and the Simplicial Volume of a Knot}, Acta Mathematica 186 (2001), 85--104; arXiv:math/9905075v2, introduction, normalization before evaluation at the root of unity, and Section 5, Conjecture 5.1. \url{https://arxiv.org/abs/math/9905075}.
\end{description}
\label{end:1485}
\end{document}
